\documentclass[11pt]{amsart}

\usepackage{amsmath,amssymb,amsthm}
\usepackage{hyperref}
\usepackage{booktabs}

\newtheorem{theorem}{Theorem}[section]
\newtheorem{lemma}[theorem]{Lemma}
\newtheorem{proposition}[theorem]{Proposition}
\newtheorem{corollary}[theorem]{Corollary}
\newtheorem{conjecture}[theorem]{Conjecture}
\theoremstyle{definition}
\newtheorem{definition}[theorem]{Definition}
\newtheorem{remark}[theorem]{Remark}
\newtheorem{example}[theorem]{Example}

\DeclareMathOperator{\CV}{CV}
\DeclareMathOperator{\Var}{Var}
\DeclareMathOperator{\arctanh}{arctanh}
\DeclareMathOperator{\arctan}{arctan}
\DeclareMathOperator{\Tr}{Tr}

\title{The Cayley Transform, Delannoy Paths, and the\\Fourier Analysis of Hamiltonian Paths in Random Tournaments}

\author{Eliott Cassidy}
\thanks{With computational assistance from Claude (Anthropic).}

\begin{document}

\begin{abstract}
We prove that the squared coefficient of variation of the Hamiltonian path count $H(T)$ over random tournaments on $n$ vertices satisfies
\[
\CV^2(H) = \sum_{k=1}^{\lfloor(n-1)/2\rfloor} \frac{[x^k]\, Q(x)^{n-2k}}{(n)_{2k}},
\]
where $Q(x) = (1+x)/(1-x)$ is the Cayley transform. We establish the master generating function
\[
Q(x)^m = 1 + 2\sum_{k=1}^{\infty} g_k(m)\, x^k,
\]
identify $k\cdot g_k(m) = \sum_j j\binom{k}{j}\binom{m}{j}2^{j-1}$ as the total diagonal-step count across all Delannoy lattice paths from $(0,0)$ to $(k,m)$ (OEIS A108666 on the diagonal), and derive a Rodrigues formula, a duality $k\,g_k(m) = m\,g_m(k)$, a parity theorem, and a functional equation $Q^m Q(-x)^m = 1$ that makes even-index weights algebraically redundant. The underlying generator is $\arctanh(x) = x + x^3/3 + x^5/5 + \cdots$, the unique odd formal power series with rational exponential. The Wick rotation $x \to ix$ sends $\arctanh$ to $i\arctan$, converting the tournament (hyperbolic, temporal, divergent) into the circle (spherical, spatial, convergent), and producing $\pi$ via $\arctanh(i) = i\pi/4$.
\end{abstract}

\maketitle

%% ====================================================================
\section{Introduction}
%% ====================================================================

A \emph{tournament} $T$ on $n$ vertices is a complete oriented graph: for each pair $\{u,v\}$, exactly one of $u \to v$ or $v \to u$ holds. R\'edei's theorem~\cite{redei1934} states that the number $H(T)$ of directed Hamiltonian paths is always odd.

When $T$ is drawn uniformly from all $2^{\binom{n}{2}}$ tournaments, $H(T)$ becomes a random variable with $\mathbb{E}[H] = n!/2^{n-1}$. This paper gives a complete description of $\CV^2(H) = \Var(H)/\mathbb{E}[H]^2$ through the lens of the Cayley transform $Q(x) = (1+x)/(1-x) = \exp(2\arctanh x)$.

\subsection{Main results}

\begin{theorem}[Master Generating Function]\label{thm:master}
$\displaystyle Q(x)^m = \left(\frac{1+x}{1-x}\right)^m = 1 + 2\sum_{k=1}^{\infty} g_k(m)\, x^k$, where $g_k(m) = \sum_{j=1}^{\min(k,m)} \binom{k-1}{j-1}\binom{m}{j}2^{j-1}$.
\end{theorem}

\begin{theorem}[Weight Formula]\label{thm:weight}
For a domino subset $S \subseteq \{0,\ldots,n-2\}$ with $|S|=L$ and $c$ connected components, $\mathbb{E}[\prod_{j \in S} Z_j] = 2^c/(n)_L$, where $Z_j = \mathbf{1}[\sigma(j{+}1)=\sigma(j){+}1] - \mathbf{1}[\sigma(j{+}1)=\sigma(j){-}1]$ for a uniform random permutation $\sigma$.
\end{theorem}

\begin{theorem}[Variance Formula]\label{thm:cv2}
$\displaystyle \CV^2(H) = \sum_{k=1}^{\lfloor(n-1)/2\rfloor} \frac{2\,g_k(n-2k)}{(n)_{2k}} = \frac{2}{n} - \frac{14}{3n^3} + O(n^{-4})$.
\end{theorem}

From these, we derive:

\begin{corollary}[Delannoy Identity]\label{cor:delannoy}
$k\,g_k(m) = \sum_{j=1}^{\min(k,m)} j\binom{k}{j}\binom{m}{j}2^{j-1}$ counts total diagonal steps in Delannoy paths from $(0,0)$ to $(k,m)$.
\end{corollary}

\begin{corollary}[Duality]\label{cor:duality}
$k\,g_k(m) = m\,g_m(k)$.
\end{corollary}

\begin{corollary}[Parity]\label{cor:parity}
$g_k(-m) = (-1)^k g_k(m)$.
\end{corollary}

\begin{corollary}[Rodrigues Formula]\label{cor:rodrigues}
$g_k(m) = \frac{1}{k!}\left[\frac{d^k}{du^k}\,u^m(u+1)^{k-1}\right]_{u=1}$.
\end{corollary}

\begin{corollary}[Functional Equation]\label{cor:functeq}
$Q(x)^m \cdot Q(-x)^m = 1$, so $g_{2k}$ is algebraically determined by $g_1, g_3, \ldots, g_{2k-1}$.
\end{corollary}

\begin{corollary}[Wick Rotation]\label{cor:wick}
$\arctanh(ix) = i\arctan(x)$. In particular, $\arctanh(i) = i\pi/4$ and $Q(i)^m = i^m$.
\end{corollary}

%% ====================================================================
\section{The $Z$-process}
%% ====================================================================

Fix $n \geq 2$ and a uniform random permutation $\sigma$ of $\{0,\ldots,n{-}1\}$. Define
\[
X_j = \mathbf{1}[\sigma(j{+}1)=\sigma(j){+}1],\quad Y_j = \mathbf{1}[\sigma(j{+}1)=\sigma(j){-}1],\quad Z_j = X_j - Y_j \in \{-1,0,1\}.
\]

Since $X_jY_j = 0$, we have $(1+Z_j) = (1+X_j)(1-Y_j)$. Define
\[
W(n) = \sum_{\sigma \in S_n} \prod_{j=0}^{n-2}(1+Z_j(\sigma)),
\]
which sums $2^{\#\text{unit ascents}}$ over permutations with no unit descents. The Fourier analysis of $H(T)$ yields $W(n)/n! = 1 + \CV^2(H)$.

\begin{definition}
A subset $S \subseteq \{0,\ldots,n{-}2\}$ is a \emph{domino subset} if every element has a neighbor in $S$.
\end{definition}

Expanding the product: $W(n)/n! = \sum_S \mathbb{E}[\prod_{j \in S} Z_j]$. By the complement symmetry $Z_j \mapsto -Z_j$, odd-size subsets contribute zero; among even-size subsets, only domino subsets contribute.

%% ====================================================================
\section{Proof of the Weight Formula}
%% ====================================================================

\begin{proof}[Proof of Theorem~\ref{thm:weight}]
\textbf{Single block.} For a contiguous block $S = \{a,\ldots,a{+}L{-}1\}$ of even size $L$, the product $\prod Z_j \neq 0$ requires each $Z_j \in \{+1,-1\}$, so $\sigma(a),\ldots,\sigma(a{+}L)$ form a $\pm 1$ walk. For distinct values (as $\sigma$ is a permutation), the walk must be monotone. Ascending and descending each give $n{-}L$ starting values with product $+1$ (since $L$ is even), and $(n{-}L{-}1)!$ arrangements of remaining values:
\[
\mathbb{E}\left[\prod_{j \in S} Z_j\right] = \frac{2(n-L)!}{n!} = \frac{2}{(n)_L}.
\]

\textbf{Multiple blocks.} For $c$ separated blocks of total size $L$, each block requires a contiguous range of integer values with $2$ orientations. The number of ways to place $c$ distinguishable disjoint value-intervals in $\{0,\ldots,n{-}1\}$ is $(n{-}L)!/(n{-}L{-}c)!$ by stars-and-bars, giving
\[
\mathbb{E}\left[\prod_{j \in S} Z_j\right] = \frac{2^c(n-L)!}{n!} = \frac{2^c}{(n)_L}.\qedhere
\]
\end{proof}

%% ====================================================================
\section{Cluster counting and the closed form}
%% ====================================================================

\begin{lemma}\label{lem:cluster}
The number of $k$-matchings of the path $P_{m+2k-2}$ with exactly $j$ clusters is $\binom{k-1}{j-1}\binom{m}{j}$.
\end{lemma}

\begin{proof}
A $j$-cluster matching is specified by a composition of $k$ into $j$ parts ($\binom{k-1}{j-1}$ choices) and a placement of $j$ clusters on the path. Each cluster of size $s$ occupies $2s$ vertex positions. With $m{-}1$ free positions distributed into $j{-}1$ mandatory gaps and $2$ margins: $\binom{m}{j}$ placements.
\end{proof}

\begin{proof}[Proof of Theorem~\ref{thm:master}]
By the weight formula, a $k$-matching with $j$ clusters contributes $2^j/(n)_{2k}$ to $\CV^2$. By Lemma~\ref{lem:cluster}, there are $\binom{k-1}{j-1}\binom{m}{j}$ such matchings. Summing:
\[
g_k(m) = \sum_{j=1}^{\min(k,m)} \binom{k-1}{j-1}\binom{m}{j}\,2^{j-1}.
\]
To derive the master generating function, the OGF of $g_k$ in $m$ is
\[
\sum_{m \geq 0} g_k(m)\,t^m = \frac{t}{(1-t)^2}\left(\frac{1+t}{1-t}\right)^{k-1} = \frac{t(1+t)^{k-1}}{(1-t)^{k+1}},
\]
obtained by substituting $\sum_{m \geq j}\binom{m}{j}t^m = t^j/(1-t)^{j+1}$ and applying the binomial theorem. The GF in $k$ follows by partial fractions:
\[
\sum_{k \geq 1} g_k(m)\,x^k = \frac{Q(x)^m - 1}{2},
\]
which gives $Q(x)^m = 1 + 2\sum_{k \geq 1} g_k(m)\,x^k$.
\end{proof}

%% ====================================================================
\section{Corollaries}
%% ====================================================================

\begin{proof}[Proof of Corollary~\ref{cor:delannoy}]
$k\binom{k-1}{j-1} = j\binom{k}{j}$.
\end{proof}

\begin{proof}[Proof of Corollary~\ref{cor:duality}]
$j\binom{k}{j}\binom{m}{j}$ is symmetric in $k$ and $m$.
\end{proof}

\begin{proof}[Proof of Corollary~\ref{cor:parity}]
$Q(-x) = (1-x)/(1+x) = Q(x)^{-1}$. So $Q(-x)^m = Q(x)^{-m}$. Comparing coefficients: $(-1)^k g_k(m) = g_k(-m)$.
\end{proof}

\begin{proof}[Proof of Corollary~\ref{cor:rodrigues}]
By Cauchy's integral formula under the substitution $u = Q(x)$, $x = (u-1)/(u+1)$: $g_k(m) = \frac{1}{k!}[d^k/du^k\,u^m(u+1)^{k-1}]_{u=1}$. The Delannoy identity follows by the Leibniz rule.
\end{proof}

\begin{proof}[Proof of Corollary~\ref{cor:functeq}]
$Q(x)^m Q(-x)^m = Q(x)^m Q(x)^{-m} = 1$. Expanding: $(1 + 2\sum g_k x^k)(1 + 2\sum (-1)^k g_k x^k) = 1$. At even order $2k$: $g_{2k} = \sum_{j=1}^{2k-1}(-1)^{j+1}g_j g_{2k-j}$.
\end{proof}

\begin{proof}[Proof of Corollary~\ref{cor:wick}]
$\arctanh(ix) = \sum_{k \geq 0}(ix)^{2k+1}/(2k+1) = i\sum_{k \geq 0}(-1)^k x^{2k+1}/(2k+1) = i\arctan(x)$. Setting $x = 1$: $\arctanh(i) = i\arctan(1) = i\pi/4$.
\end{proof}

%% ====================================================================
\section{The generator $\arctanh$ and the $1/n^2$ cancellation}
%% ====================================================================

\begin{proposition}[Log-coefficients]\label{prop:log}
$[x^k]\log Q(x)^m = 2m/k$ if $k$ is odd, and $0$ if $k$ is even.
\end{proposition}

\begin{proof}
$\log Q^m = 2m\arctanh(x) = 2m\sum_{k \text{ odd}} x^k/k$.
\end{proof}

The vanishing of even log-coefficients has two consequences.

\begin{corollary}[$1/n^2$ cancellation]\label{cor:cancel}
In the expansion $\CV^2 = \sum_{p \geq 1} a_p/n^p$, we have $a_2 = 0$.
\end{corollary}

\begin{proof}
The $k{=}1$ contribution is $2(n{-}2)/(n(n{-}1)) = 2/n - 2/n^2 + O(n^{-3})$. The $k{=}2$ contribution is $2(n{-}4)^2/(n)_4 = 2/n^2 + O(n^{-3})$. The cancellation $-2/n^2 + 2/n^2 = 0$ holds because $g_2(m) = g_1(m)^2 = m^2$, which follows from $[x^2]\log Q^m = 0$ (the second log-coefficient vanishes).
\end{proof}

\begin{remark}
The identity $g_2 = g_1^2$ is the simplest case of the functional equation. More generally, $g_4 = 2g_1 g_3 - g_2^2$, $g_6 = 2g_1 g_5 - 2g_2 g_4 + g_3^2$, etc. Each even-index weight is algebraically determined by odd-index weights, which are themselves determined by $\arctanh$.
\end{remark}

%% ====================================================================
\section{The transfer matrix and the $x$-tribonacci recurrence}
%% ====================================================================

\begin{proposition}\label{prop:transfer}
The generating function $F(N,x) = \sum_k [x^k]Q(x)^{N-2k+2}\cdot x^k$ satisfies
\[
F(N{+}3,x) = F(N{+}2,x) + x\,F(N{+}1,x) + x\,F(N,x)
\]
with generating function $\sum_{N \geq 0} F(N,x)\,y^N = (1+2xy+xy^2)/(1-y-xy^2-xy^3)$.
\end{proposition}

This follows from the transfer matrix $M(x) = \left(\begin{smallmatrix}1&2x&0\\0&0&1\\1&x&0\end{smallmatrix}\right)$, whose characteristic polynomial $\lambda^3 - \lambda^2 - x\lambda - x = 0$ gives the denominator $1 - y - xy^2 - xy^3 = \det(I - yM)$. At $x = 1$, the dominant eigenvalue is the tribonacci constant $\tau = 1.8393\ldots$, satisfying $\tau^3 = \tau^2 + \tau + 1$.

The three eigenvalues are:
\begin{itemize}
\item $\lambda_1(x) = 1 + 2x + O(x^2)$: real, dominant (Perron--Frobenius);
\item $\lambda_{2,3}(x) \sim \pm i\sqrt{x}$: complex conjugate pair.
\end{itemize}

The complex pair arises from the degeneracy splitting of the null state at $x = 0$. The tribonacci constant $\tau = \varphi + 0.2213\ldots$ where $\varphi = (1+\sqrt{5})/2$ is the golden ratio; the correction $\tau - \varphi$ is the \emph{cost of memory}.

%% ====================================================================
\section{The Wick rotation and the emergence of $\pi$}
%% ====================================================================

The functions $\arctanh$ and $\arctan$ share the same coefficients but differ by sign:
\begin{align*}
\arctanh(x) &= x + \frac{x^3}{3} + \frac{x^5}{5} + \cdots && \text{(constant signs: \emph{temporal})}, \\
\arctan(x) &= x - \frac{x^3}{3} + \frac{x^5}{5} - \cdots && \text{(alternating signs: \emph{spatial})}.
\end{align*}

The Wick rotation $x \to ix$ converts one to the other: $\arctanh(ix) = i\arctan(x)$.

\begin{proposition}\label{prop:pi}
$Q(i) = i$ and $\arctanh(i) = i\pi/4$. More generally, $Q(i)^m = i^m = e^{im\pi/2}$, so the period-4 orbit $\{1, i, -1, -i\}$ of $Q$ in $\mathrm{PGL}(2,\mathbb{C})$ parametrizes the unit circle by $m\pi/2$.
\end{proposition}

\begin{proof}
$Q(i) = (1+i)/(1-i) = i$. Then $\exp(2\arctanh(i)) = i$, so $\arctanh(i) = \frac{1}{2}\log i = i\pi/4$.
\end{proof}

This provides a dictionary:
\begin{center}
\renewcommand{\arraystretch}{1.3}
\begin{tabular}{@{}lll@{}}
\toprule
& \textbf{Tournament (real $x$)} & \textbf{Circle (imaginary $ix$)} \\
\midrule
Generator & $\arctanh(x)$ & $i\arctan(x)$ \\
Signs & all positive & alternating \\
Geometry & hyperbolic & spherical \\
Convergence & diverges at $x=1$ & converges to $\pi/4$ \\
Interpretation & temporal (direction) & spatial (symmetry) \\
\bottomrule
\end{tabular}
\end{center}

The constant signs in $\arctanh$ produce divergence (infinite time: one can never finish comparing). The alternating signs in $\arctan$ produce convergence (finite space: the circle closes). The number $\pi$ emerges from the tournament at imaginary coupling.

%% ====================================================================
\section{The sequence $W(n)$ and connections}
%% ====================================================================

The sequence $W(n) = n!(1 + \CV^2)$ begins $1, 2, 8, 32, 158, 928, 6350, 49752, 439670, \ldots$ and does not appear in the OEIS as of 2026.

The diagonal $T(k,k) = k\,g_k(k)$ is OEIS A108666 (total diagonal steps in Delannoy paths of length $k$).

\subsection{Summary of formulas}

\begin{center}
\renewcommand{\arraystretch}{1.4}
\begin{tabular}{@{}cl@{}}
\toprule
\# & Formula \\
\midrule
1 & $Q(x)^m = 1 + 2\sum g_k(m)\,x^k$, \quad $Q = (1{+}x)/(1{-}x) = e^{2\arctanh x}$ \\
2 & $k\,g_k(m) = \sum j\binom{k}{j}\binom{m}{j}2^{j-1}$ \quad (Delannoy diagonal steps) \\
3 & $g_k(m) = \frac{1}{k!}[d^k/du^k\,u^m(u{+}1)^{k-1}]_{u=1}$ \quad (Rodrigues) \\
4 & $\sum g_k(m)\,t^m = t(1{+}t)^{k-1}/(1{-}t)^{k+1}$ \quad (OGF in $m$) \\
5 & $Q^m Q(-x)^m = 1$ \quad (functional equation, even $g_k$ redundant) \\
6 & $[x^k]\log Q^m = 2m/k$ for odd $k$, $0$ for even $k$ \\
7 & $k\,g_k(m) = m\,g_m(k)$ \quad (duality) \\
8 & $g_k(-m) = (-1)^k g_k(m)$ \quad (parity) \\
9 & $\arctanh(ix) = i\arctan(x)$; \quad $\arctanh(i) = i\pi/4$ \quad (Wick rotation) \\
10 & $\CV^2 = 2/n + 0/n^2 - 14/(3n^3) + O(n^{-4})$ \\
\bottomrule
\end{tabular}
\end{center}

\subsection{Uniqueness of $\arctanh$}

\begin{proposition}\label{prop:unique}
$\arctanh$ is the unique odd formal power series whose exponential is rational of degree $(1,1)$.
\end{proposition}

\begin{proof}
If $f(-x) = -f(x)$ and $e^f = P/Q$ with $P, Q$ coprime polynomials, then $e^{-f} = e^{f(-x)} = P(-x)/Q(-x)$, so $P(x)Q(-x) = P(-x)Q(x)$. By coprimality, $P(-x) = cQ(x)$, and $c = 1$ from $f(0) = 0$. For degree $(1,1)$: $Q = 1 - ax$, $P = 1 + ax$, $f = \log((1{+}ax)/(1{-}ax)) = 2\arctanh(ax)$.
\end{proof}

%% ====================================================================
\section{Interpretation}
%% ====================================================================

The entire tournament Fourier theory reduces to the sequence $1, 1/3, 1/5, 1/7, \ldots$ of reciprocal odd integers---the Taylor coefficients of $\arctanh$. These encode the \emph{directed interaction scales}: the weight $1/(2k{+}1)$ at scale $2k{+}1$ measures the contribution of directed cycles of that length.

The even coefficients of $\log Q^m$ vanish because the tournament is built from independent arc decisions, and even-order interactions are symmetric (insensitive to direction). The directed (odd) and undirected (even) parts decompose the theory into
\[
Q^m = \underbrace{\cosh(2m\arctanh x)}_{\text{symmetric (space)}} + \underbrace{\sinh(2m\arctanh x)}_{\text{antisymmetric (time)}},
\]
and the Wick rotation $x \to ix$ converts the temporal $\arctanh$ into the spatial $\arctan$, producing $\pi$.

%% ====================================================================
\section{The Golden Shadow}
%% ====================================================================

The metallic means $y_n = (n + \sqrt{n^2+4})/2$ (with $y_1 = \varphi$, $y_2 = 1+\sqrt{2}$, $y_3 = (3+\sqrt{13})/2$, etc.)\ arise throughout number theory and tiling. We define their \emph{shifted companions}:

\begin{definition}\label{def:fn}
For each integer $n \geq 1$, define
\[
f_n \;=\; \frac{n - 2 + \sqrt{n^2 + 4}}{2} \;=\; y_n - 1.
\]
\end{definition}

\begin{proposition}[Quadratic and continued fraction]\label{prop:fn-quad}
$f_n$ satisfies the quadratic $f^2 - (n-2)f - n = 0$, and has the purely periodic continued fraction
\[
f_n = [n-1;\, \overline{n, n, n, \ldots}].
\]
\end{proposition}

\begin{proof}
Since $y_n$ satisfies $y^2 - ny - 1 = 0$, the substitution $f = y - 1$ gives
\[
(f+1)^2 - n(f+1) - 1 = f^2 + 2f + 1 - nf - n - 1 = f^2 - (n-2)f - n = 0.
\]
For the continued fraction, note that $f_n = n - 1 + 1/f_n'$ where $f_n' = n + 1/f_n'$, giving $f_n' = [n; \overline{n}] = (\text{purely periodic})$. Combining: $f_n = [n-1; n, n, n, \ldots]$.
\end{proof}

\begin{proposition}[The Cayley collapse]\label{prop:cayley-collapse}
For each $f = f_n$, the rational expression $(f^2 + f - 1)/(f^2 + 3f + 1)$ collapses to a rational number:
\[
\frac{f^2 + f - 1}{f^2 + 3f + 1} \;=\; \frac{n-1}{n+1}.
\]
The irrational part $\sqrt{n^2+4}$ cancels because the common factor $(f+1)$ divides both numerator and denominator.
\end{proposition}

\begin{proof}
Using $f^2 = (n-2)f + n$ (from Proposition~\ref{prop:fn-quad}):
\begin{align*}
f^2 + f - 1 &= (n-2)f + n + f - 1 = (n-1)f + (n-1) = (n-1)(f+1), \\
f^2 + 3f + 1 &= (n-2)f + n + 3f + 1 = (n+1)f + (n+1) = (n+1)(f+1).
\end{align*}
Therefore $(f^2 + f - 1)/(f^2 + 3f + 1) = (n-1)(f+1)/((n+1)(f+1)) = (n-1)/(n+1)$.
\end{proof}

\begin{example}[Special cases]\label{ex:fn-special}
\leavevmode
\begin{itemize}
\item $f_1 = (-1 + \sqrt{5})/2 = 1/\varphi$, where $\varphi = (1+\sqrt{5})/2$ is the golden ratio.
\item $f_2 = \sqrt{2}$, the silver companion. Here $f_2^2 = 2$, confirming $f^2 - 0 \cdot f - 2 = 0$.
\item $f_4 = (2 + \sqrt{20})/2 = 1 + \sqrt{5} = \varphi^2$.
\end{itemize}
\end{example}

%% ====================================================================
\section{Exceptional Points and the Golden Ratio}
%% ====================================================================

Returning to the transfer matrix $M(x)$ of Section~7, the characteristic polynomial is $p(\lambda) = \lambda^3 - \lambda^2 - x\lambda - x$. Its discriminant with respect to $\lambda$ is:

\begin{theorem}[Discriminant and exceptional points]\label{thm:disc}
\[
\Delta(x) \;=\; 4x(x^2 - 11x - 1).
\]
The transfer matrix $M(x)$ has a repeated eigenvalue (an \emph{exceptional point}) at exactly three values of $x$:
\begin{enumerate}
\item $x = 0$: the trivial degeneracy (two eigenvalues collide at $\lambda = 0$), with degenerate eigenvalue $d = 0$;
\item $x = x_- = \frac{11 - 5\sqrt{5}}{2} = 8 - 5\varphi$: with degenerate eigenvalue $d = 1/\varphi$;
\item $x = x_+ = \frac{11 + 5\sqrt{5}}{2} = 3 + 5\varphi$: with degenerate eigenvalue $d = -\varphi$.
\end{enumerate}
\end{theorem}

\begin{proof}
The discriminant of $\lambda^3 + a\lambda^2 + b\lambda + c$ is $\Delta = 18abc - 4a^3c + a^2b^2 - 4b^3 - 27c^2$. With $a = -1$, $b = -x$, $c = -x$:
\begin{align*}
\Delta &= 18(-1)(-x)(-x) - 4(-1)^3(-x) + (-1)^2(-x)^2 - 4(-x)^3 - 27(-x)^2 \\
       &= -18x^2 - 4x + x^2 + 4x^3 - 27x^2 \\
       &= 4x^3 - 44x^2 - 4x \\
       &= 4x(x^2 - 11x - 1).
\end{align*}
The nontrivial exceptional points are the roots of $x^2 - 11x - 1 = 0$, namely $x_\pm = (11 \pm \sqrt{125})/2 = (11 \pm 5\sqrt{5})/2$. At each EP, the repeated eigenvalue $d$ satisfies $3d^2 - 2d - x = 0$ (from $p'(d) = 0$) and $d^3 - d^2 - xd - x = 0$ (from $p(d) = 0$). Eliminating $x = 3d^2 - 2d$ from $p(d) = 0$:
\[
d^3 - d^2 - d(3d^2 - 2d) - (3d^2 - 2d) = 0 \implies -2d^3 - 2d^2 + 2d = 0 \implies d(d^2 + d - 1) = 0.
\]
\end{proof}

\begin{remark}[The golden ratio equation]\label{rem:golden}
The EP eigenvalue equation $d^2 + d - 1 = 0$ has roots $d = (-1 \pm \sqrt{5})/2$, which are $1/\varphi$ and $-\varphi$. This is precisely the \emph{golden ratio equation}: the same algebraic relation that defines $\varphi$ governs the coalescence of eigenvalues in the tournament transfer matrix. One matrix simultaneously encodes:
\begin{itemize}
\item the tribonacci constant $\tau$ (at $x = 1$, real-world tournament counting),
\item the golden ratio $\varphi$ (at the exceptional points, where eigenvalue degeneracy occurs),
\item exceptional-point physics (non-Hermitian spectral theory, $\mathcal{PT}$-symmetry breaking).
\end{itemize}
\end{remark}

\begin{remark}[The $-44$ coincidence]\label{rem:minus44}
At $x \in \{-1, 1, 11\}$, the discriminant evaluates to
\[
\Delta(-1) = \Delta(1) = \Delta(11) = -44.
\]
These three points are: $x = -1$ (the zero of $Q$, where $Q(-1) = 0$), $x = 1$ (the pole of $Q$, where $Q(1) = \infty$, and the tribonacci specialization), and $x = 11$ (the Paley prime $p = 11$, the smallest prime where the Paley tournament exhibits non-trivial structure in this theory). The coincidence $\Delta = -44$ at all three values follows from $x^2 - 11x - 1$ evaluating to $11, -11, 9$ respectively, combined with the prefactor $4x$.
\end{remark}

%% ====================================================================
\section{Connections}
%% ====================================================================

The Cayley transform $Q(x) = (1+x)/(1-x)$ and the underlying $\arctanh$ generator connect the tournament theory to several independent areas.

\subsection{The bilinear (Tustin) transform}

In digital signal processing, the \emph{bilinear transform} (also called the Tustin transform) maps an analog transfer function $H_a(s)$ to its digital counterpart $H_d(z)$ via the substitution
\[
s \;=\; \frac{2}{T}\cdot\frac{z - 1}{z + 1},
\]
where $T$ is the sampling period. Up to the scaling factor $2/T$, the map $z \mapsto (z-1)/(z+1)$ is exactly the \emph{inverse Cayley transform} $Q^{-1}(z) = (z-1)/(z+1)$, and the forward map is $z = Q(s) = (1+s)/(1-s)$ (with $T = 2$).

This identification appears to be genuinely new: while both objects are classical, no existing literature (to the authors' knowledge) connects the tournament Fourier analysis of Sections~2--6 with the bilinear/Tustin transform of DSP theory. The Cayley transform $Q$ serves as the bridge: its role as the moment generating function of the $Z$-process (tournament side) is the same M\"obius transformation that maps the analog $s$-plane to the digital $z$-plane (signal processing side). The frequency warping $\omega_d = \arctan(\omega_a T/2)$ in DSP is precisely the Wick rotation of Corollary~\ref{cor:wick}.

\subsection{Information geometry: Fisher--Rao metric}

The Fisher information of the Bernoulli distribution $\mathrm{Bern}(p)$ is $I(p) = 1/(p(1-p))$, and the induced Fisher--Rao geodesic distance between $p$ and $q$ is
\[
d_{\mathrm{FR}}(p,q) \;=\; 2\left|\arctanh(2p - 1) - \arctanh(2q - 1)\right|.
\]
At the symmetric point $p = 1/2$ (the uniform tournament measure), the curvature of the statistical manifold equals $\CV^2$. The Fisher--Rao metric is thus the natural metric for the tournament parameter space, with $\arctanh$ as the coordinate map.

\subsection{Universality class}

The following objects all share the same underlying algebraic structure---powers of the M\"obius transformation $Q(x) = (1+x)/(1-x)$, its logarithm $\arctanh$, or its Wick rotation $\arctan$:
\begin{enumerate}
\item Tournament $\CV^2$ (this paper);
\item Central Delannoy numbers and lattice paths~\cite{sulanke, hetyei2009};
\item The bilinear/Tustin transform in DSP~\cite{marcus-roth-siegel};
\item Fisher--Rao distance on $\mathrm{Bern}(p)$;
\item The Gauss continued fraction $\arctanh(1/n) = [0; n, 3n, 5n, \ldots]$;
\item Legendre polynomials via the Rodrigues formula (Corollary~\ref{cor:rodrigues} at $u = 1$);
\item Hyperbolic geometry (half-plane model, cross-ratio);
\item The Leibniz--Gregory--Madhava series $\pi/4 = 1 - 1/3 + 1/5 - \cdots$.
\end{enumerate}
That these eight domains are governed by a single function is not a metaphor: they share the same coefficients, the same recurrences, and the same functional equation.

\begin{thebibliography}{99}
\bibitem{redei1934} L.~R\'edei, \emph{Ein kombinatorischer Satz}, Acta Litt.\ Szeged \textbf{7} (1934), 39--43.

\bibitem{grinberg-stanley} D.~Grinberg and R.~P.~Stanley, \emph{On the number of Hamiltonian paths in a tournament}, arXiv:2307.05569 (2023).

\bibitem{oeis-a108666} OEIS Foundation, \emph{A108666: Number of $(1,1)$-steps in all Delannoy paths of length $n$}, \url{https://oeis.org/A108666}.

\bibitem{sulanke} R.~A.~Sulanke, \emph{Objects counted by the central Delannoy numbers}, J.\ Integer Seq.\ \textbf{6} (2003), Art.~03.1.5.

\bibitem{hetyei2009} G.~Hetyei, \emph{Central Delannoy numbers, Legendre polynomials, and a balanced join operation preserving the Cohen--Macaulay property}, Ann.\ Comb.\ \textbf{13} (2009), 363--396.

\bibitem{marcus-roth-siegel} B.~H.~Marcus, R.~M.~Roth, and P.~H.~Siegel, \emph{Constrained systems and coding for recording channels}, in \emph{Handbook of Coding Theory} (V.~S.~Pless and W.~C.~Huffman, eds.), vol.~2, Elsevier, 1998, pp.~1635--1764.
\end{thebibliography}

\end{document}
